

\ch{Appendix C: Acronyms and Terms}\label{appendix:acronyms-terms}

This appendix provides two references for the reader. The first is a list of acronyms and abbreviations used throughout the dissertation. The second is a glossary of disciplinary terms, each accompanied by a concise definition intended to make the technical narrative accessible to readers across systems engineering, Information Assurance, computer science, and IT service management. Established artifacts and frameworks retain the names under which they were published; consistent with the convention noted in Chapter~1, the institution itself appears as the Department of War (DoW).

\section{Acronyms}

\begin{singlespace}
\begin{longtable}{p{0.16\textwidth} p{0.76\textwidth}}
\toprule
\textbf{Acronym} & \textbf{Definition} \\
\midrule
\endfirsthead
\multicolumn{2}{c}{\textit{(Acronyms, continued from previous page)}} \\
\toprule
\textbf{Acronym} & \textbf{Definition} \\
\midrule
\endhead
\midrule
\multicolumn{2}{r}{\textit{continued on next page}} \\
\endfoot
\bottomrule
\endlastfoot
ACM       & Association for Computing Machinery \\
ANOVA     & Analysis of Variance \\
API       & Application Programming Interface \\
AVE       & Average Variance Extracted \\
BoE       & Body of Evidence \\
CAE       & Center of Academic Excellence \\
CFI       & Comparative Fit Index \\
CFR       & Code of Federal Regulations \\
CISA      & Cybersecurity and Infrastructure Security Agency \\
CM        & Configuration Management \\
CMDB      & Configuration Management Database \\
CMMC      & Cybersecurity Maturity Model Certification \\
COBIT     & Control Objectives for Information and Related Technologies \\
CPS       & Cyber-Physical System \\
CR        & Composite Reliability \\
CSA       & Cloud Security Alliance \\
CSF       & Cybersecurity Framework (NIST) \\
CUI       & Controlled Unclassified Information \\
DAU       & Defense Acquisition University \\
DE        & Digital Engineering \\
DevOps    & Development and Operations \\
DevSecOps & Development, Security, and Operations \\
DIF       & Differential Item Functioning \\
DISA      & Defense Information Systems Agency \\
DMM       & Domain Metamodel (UAF) \\
DoD       & Department of Defense \\
DoDAF     & Department of Defense Architecture Framework \\
DoDI      & Department of Defense Instruction \\
DoW       & Department of War \\
DSU       & Dakota State University \\
EA        & Enterprise Architecture \\
FAR       & Federal Acquisition Regulation \\
FEAF      & Federal Enterprise Architecture Framework \\
FERPA     & Family Educational Rights and Privacy Act \\
GDPR      & General Data Protection Regulation \\
GRC       & Governance, Risk, and Compliance \\
GSN       & Goal Structuring Notation \\
HC3       & Heteroscedasticity-Consistent standard errors (type 3) \\
HIPAA     & Health Insurance Portability and Accountability Act \\
HTMT      & Heterotrait-Monotrait ratio (of correlations) \\
IA        & Information Assurance \\
IEC       & International Electrotechnical Commission \\
IEEE      & Institute of Electrical and Electronics Engineers \\
IETF      & Internet Engineering Task Force \\
INCOSE    & International Council on Systems Engineering \\
IoT       & Internet of Things \\
IRB       & Institutional Review Board \\
ISACA     & Information Systems Audit and Control Association \\
ISO       & International Organization for Standardization \\
IT        & Information Technology \\
ITIL      & Information Technology Infrastructure Library \\
ITSM      & Information Technology Service Management \\
MBSE      & Model-Based Systems Engineering \\
ML        & Modeling Language (UAF usage; not machine learning) \\
MODAF     & Ministry of Defence Architecture Framework (United Kingdom) \\
NAF       & NATO Architecture Framework \\
NASA      & National Aeronautics and Space Administration \\
NATO      & North Atlantic Treaty Organization \\
NDIA      & National Defense Industrial Association \\
NIST      & National Institute of Standards and Technology \\
NSA       & National Security Agency \\
NSTISSI   & National Security Telecommunications and Information Systems Security Instruction \\
OMG       & Object Management Group \\
OSCAL     & Open Security Controls Assessment Language \\
OSF       & Open Science Framework \\
PCI DSS   & Payment Card Industry Data Security Standard \\
PLM       & Product Lifecycle Management \\
PLS-SEM   & Partial Least Squares Structural Equation Modeling \\
RMF       & Risk Management Framework \\
RMSEA     & Root Mean Square Error of Approximation \\
ROI       & Return on Investment \\
RQ        & Research Question \\
SACM      & Structured Assurance Case Metamodel \\
SD        & Standard Deviation \\
SE        & Systems Engineering \\
SEBoK     & Systems Engineering Body of Knowledge \\
SEI       & Software Engineering Institute \\
SERC      & Systems Engineering Research Center \\
SoS       & System of Systems \\
SP        & Special Publication (NIST) \\
SSP       & System Security Plan \\
SysML     & Systems Modeling Language \\
TAM       & Technology Acceptance Model \\
TOGAF     & The Open Group Architecture Framework \\
UAF       & Unified Architecture Framework \\
ULMC      & Unmeasured Latent Method Construct \\
UML       & Unified Modeling Language \\
UPDM      & Unified Profile for DoDAF and MODAF \\
UTAUT     & Unified Theory of Acceptance and Use of Technology \\
\end{longtable}
\end{singlespace}

\section{Terms and Definitions}

\begin{singlespace}
\begin{longtable}{p{0.26\textwidth} p{0.66\textwidth}}
\toprule
\textbf{Term} & \textbf{Definition} \\
\midrule
\endfirsthead
\multicolumn{2}{c}{\textit{(Terms, continued from previous page)}} \\
\toprule
\textbf{Term} & \textbf{Definition} \\
\midrule
\endhead
\midrule
\multicolumn{2}{r}{\textit{continued on next page}} \\
\endfoot
\bottomrule
\endlastfoot
Acquiescence bias & The tendency of survey respondents to agree with statements regardless of content, which can inflate agreement on affirmatively framed items. \\
Assurance case & A structured, evidence-supported argument that a system satisfies a claimed property such as safety or security. \\
Authoritative source of truth & A single, trusted baseline of system data and models from which all stakeholders operate, ensuring consistency and traceability across the lifecycle. \\
Authorization boundary & The set of information system components covered by a security authorization, defining the scope of risk-management responsibility under the Risk Management Framework. \\
Boundary object & An artifact, such as a shared model or specification, flexible enough to serve local needs yet stable enough to coordinate work across different communities of practice. \\
Common-method variance & The portion of observed covariance among measures attributable to the shared self-report method rather than to the constructs, screened here with Harman's single-factor test. \\
Composite score & A summary measure formed by averaging a set of related Likert items; this study uses the Awareness, IT Value, and Security Value composites. \\
Configuration management & The discipline of establishing and maintaining the integrity of a system's components, attributes, and relationships through controlled change. \\
Continuous monitoring & The ongoing assessment of security-control effectiveness and system state required by the Risk Management Framework. \\
Convenience sampling & A non-probability sampling method that recruits accessible participants, which constrains the generalizability of results. \\
Cram\'er's V & An effect-size measure of association for a chi-square test of independence, ranging from 0 to 1. \\
Cronbach's alpha & A coefficient of internal-consistency reliability for a set of items, with values of 0.70 or higher conventionally considered acceptable. \\
Differential item functioning & A condition in which respondents from different groups but at the same underlying trait level answer an item differently, screened here across professional field. \\
Digital Engineering & An integrated, model-based approach to designing, developing, and sustaining complex systems through authoritative digital models and data. \\
Digital thread & The authoritative, traceable connections among requirements, design, implementation, and validation artifacts that persist across a system's lifecycle. \\
Digital twin & A virtual representation of a physical or logical system that maintains synchronization with its real-world counterpart through continuous data exchange. \\
Effect size & A measure of the magnitude of a difference or association, such as Cohen's $d$ or Cram\'er's $V$, that is independent of sample size. \\
Enterprise architecture & A discipline and set of artifacts describing the structure and relationships among an organization's business processes, information, applications, and technology. \\
Eta-squared ($\eta^2$) & A non-parametric effect-size measure for the Kruskal--Wallis test, expressing the share of rank variance attributable to group membership. \\
Holm--Bonferroni correction & A sequential procedure that adjusts significance thresholds to control the family-wise error rate across multiple related comparisons. \\
Information Assurance & The broader discipline of protecting and governing information and information systems, encompassing security policy, risk management, and compliance, within which cybersecurity is the operational component. \\
Kruskal--Wallis test & A non-parametric test for differences among three or more groups, used here for experience-level comparisons. \\
Likert scale & An ordered response format, here five points, used to measure attitudes, agreement, or familiarity. \\
Mann--Whitney $U$ test & A non-parametric test comparing the distributions of two independent groups without assuming normality, used here as the primary field comparison. \\
Model-Based Systems Engineering & The formalized application of digital models, rather than documents, as the primary means of capturing system requirements, design, analysis, and verification. \\
Non-response bias & Systematic error arising when survey non-respondents differ from respondents; assessed here through wave analysis. \\
Priming & In survey methodology, the influence of question wording or context on responses, diagnosed here by comparing value perceptions across familiarity strata. \\
Product Lifecycle Management & Frameworks and tools for managing a system's information, configuration baselines, and changes from conception through retirement. \\
Rank-biserial correlation & A non-parametric effect-size measure derived from the Mann--Whitney statistic. \\
Risk Management Framework & The NIST seven-step process (prepare, categorize, select, implement, assess, authorize, and monitor) for managing security and privacy risk. \\
Security control & A safeguard or countermeasure, for example from NIST SP~800-53, selected to protect the confidentiality, integrity, and availability of a system. \\
Straight-lining & A response pattern in which a respondent gives the same answer across many items; used here as an acquiescence diagnostic on the agreement-scale battery. \\
System of systems & A set of independently useful systems integrated to deliver capabilities that none could provide alone. \\
Systems engineering & An interdisciplinary engineering approach to designing, integrating, and managing complex systems across their lifecycle. \\
Unified Architecture Framework & An ISO-standardized enterprise and system-of-systems architecture framework, expressed through UML and SysML profiles, that consolidates earlier military frameworks. \\
Wave analysis & A non-response diagnostic that compares early and late survey respondents, treating late responders as proxies for non-respondents. \\
Wilson score interval & A confidence interval for a proportion that remains well-behaved near zero or one, used here for agreement rates. \\
\end{longtable}
\end{singlespace}
